\section{Square-bond exponent consequences}
\label{sec:percolation-exponents}

We prove Corollary~\ref{intro:bond-exponents} by extracting two kinds of
arm information from the complete unoriented loop limit.  With fixed
buffers at the annular rims, the absence of a surrounding loop detects
a radial arm of either color; self-duality
then compares this with one primal-open arm.  Even alternating arms
are instead detected by the number of loop traversals between the
annular rims.  These two descriptions give bounds at each fixed
annulus ratio.  Quasi-multiplicativity converts those bounds into
lattice arm exponents, after which the near-critical scaling relations
give the thermodynamic exponents.

The square model in this section has $q=1$; all critical arm
probabilities are at $p=1/2$.  On mesh $\delta$, let
$\pi^\square_{1,\delta}(a,b)$ be the
probability of an actual primal-open path from $\overline B_a$ to
$\C\setminus B_b$.  For even $m$, let $A^X_{m,\delta}(a,b)$ be the
event of $m$ disjoint cyclically alternating arms between these sets,
and write $\pi^X_{m,\delta}(a,b)$ for its probability.  Here
$X=\square$ means primal-open/dual-open square bonds, and $X=\triangle$
means the two colors of critical triangular-site percolation.
Write $\Gamma^X_\delta$ for the corresponding complete interface-loop
collection.  The annuli in this notation are circular.

\subsection{One arm from surrounding loops}
Let $B_\delta(a,b)$ be the square-bond event of a radial arm of
\emph{either} color.  For a complete loop collection $\Gamma$, let
$H_{\rm o}(\Gamma;a,b)$ be the event that some loop is contained in
$\{a<|z|<b\}$ and has nonzero absolute winding about $0$; define
$H_{\rm c}(\Gamma;a,b)$ using $\{a\le |z|\le b\}$.  Put
$Z(\Gamma;a,b)=H_{\rm o}(\Gamma;a,b)^c$, and use subscripts $\delta$
and $\mathrm{CLE}$ for the square-bond and whole-plane $\CLE_6$
collections.  Absolute winding is independent of loop orientation.

For fixed $0<a<b/4$ and all sufficiently small $\delta$, the rounded
disjoint tile loops used in Section~\ref{sec:loops} give
\begin{equation}\label{perc:one-lattice}
 B_\delta(a/2,2b)\subset Z_\delta(a,b)
 \subset B_\delta(2a,b/2).
\end{equation}
Indeed, a complete surrounding loop separates the two rims, and actual
open paths of either color stay in the complementary territories, so
neither can cross that loop.  Conversely, if neither color crosses the
cropped annulus, planar cutset duality supplies a primal-open circuit
and a dual-open circuit surrounding $0$ between its rims, up to
$O(\delta)$ displacement.  They are disjoint and nested.  The cluster
of the inner circuit's color is confined by the outer opposite-color
circuit, so its outer medial boundary is a complete surrounding loop
between them.  The fixed collars absorb the displacement.

On the locally finite matching space, $H_{\rm o}(\,\cdot\,;a,b)$ is
open: a witnessing loop has positive rim clearance and positive
spherical diameter, and a sufficiently close matched loop preserves
its absolute winding.  The event $H_{\rm c}(\,\cdot\,;a,b)$ is closed.
A closed-annulus witness has spherical diameter bounded below in terms
of $a,b$; local finiteness therefore lets a subsequence of witnesses
be matched to one fixed limit loop, which retains closed-annulus
containment and nonzero winding.  Since
$H_{\rm o}(a,b)\subset H_{\rm c}(a,b)\subset H_{\rm o}(a/2,2b)$,
Portmanteau and Theorem~\ref{intro:main}(ii) give
\begin{equation}\label{perc:one-portmanteau}
 \mathbb P(Z_{\mathrm{CLE}}(a/2,2b))
 \le \liminf_{\delta\downarrow0}\mathbb P(Z_\delta(a,b))
 \le \limsup_{\delta\downarrow0}\mathbb P(Z_\delta(a,b))
 \le \mathbb P(Z_{\mathrm{CLE}}(a,b)).
\end{equation}
This uses the open complement of $H_{\rm c}$ and the closed event $Z$;
it does not require continuity of an event indicator.  Combining
\eqref{perc:one-lattice}--\eqref{perc:one-portmanteau} at rescaled
radii yields, for $16a<b$,
\begin{equation}\label{perc:either-envelope}
\begin{split}
 \mathbb P(Z_{\mathrm{CLE}}(a/4,4b))
 &\le \liminf_{\delta\downarrow0}\mathbb P(B_\delta(a,b))\\
 &\le \limsup_{\delta\downarrow0}\mathbb P(B_\delta(a,b))
 \le \mathbb P(Z_{\mathrm{CLE}}(2a,b/2)).
\end{split}
\end{equation}

By scale invariance put $p_0(t)=\mathbb P(Z_{\mathrm{CLE}}(t,1))$.
The Camia--Newman identification of the no-surrounding-loop event with
the triangular-site either-color annular crossing limit, together with
the one-color estimate of Lawler--Schramm--Werner, gives
\begin{equation}\label{perc:cle-one}
 p_0(t)\asymp t^{5/48}\qquad(t\downarrow0).
\end{equation}
Here either-color probability lies between one-color probability and
twice that probability; fixed rim dilations accommodate either open
or closed rim conventions, as in
\citet[Theorem~2 and Section~1.1]{CamiaNewman2006Full} and
\citet[Theorem~1.2 and Section~3]{LSW2002OneArm}.
At square-bond $p=1/2$, the dual law translated by
$(-\delta/2,-\delta/2)$ is the primal law.  The shift is absorbed by
the collars, so
$\mathbb P(B_\delta(a/2,2b))/2\le\pi^\square_{1,\delta}(a,b)
\le\mathbb P(B_\delta(a,b))$ for small $\delta$.  Consequently, for
fixed $t=a/b<1/16$,
\begin{equation}\label{perc:one-envelope}
 \tfrac12p_0(t/64)
 \le\liminf_{\delta\downarrow0}\pi^\square_{1,\delta}(a,b)
 \le\limsup_{\delta\downarrow0}\pi^\square_{1,\delta}(a,b)
 \le p_0(4t).
\end{equation}

\subsection{Even alternating arms from loop traversals}
For even $m$, let $T_m(\Gamma;a,b)$ mean that there are $m$ parameter
intervals on members of $\Gamma$, with disjoint interiors on each
individual loop, each visiting both $\{|z|<a\}$ and $\{|z|>b\}$.
Intervals on distinct loops are counted separately.  This definition
uses the retained loop traversals, not merely their traces, and is
independent of roots and orientations.

For either lattice, fixed $0<a<b$, and all sufficiently small meshes,
\begin{equation}\label{perc:even-lattice}
 A^X_{m,\delta}(a/2,2b)
 \subset T_m(\Gamma^X_\delta;a,b),\qquad
 T_m(\Gamma^X_\delta;a/2,2b)
 \subset A^X_{m,\delta}(a,b).
\end{equation}
For the first inclusion, take simple last-inner-to-first-outer
subpaths of the alternating arms within the two collars.  They cut
the annulus into $m$ strips with opposite-colored
sides.  Planar separation in each strip supplies an interface
crosscut, belonging to an interval of a complete loop.  The crosscuts
are disjoint because the open arms separate the strips.  Conversely,
crop the loop intervals to disjoint radial crosscuts.  The territorial
arm counting of Lemma~\ref{chordal:territories} gives at least $m$
cyclically alternating actual arms in the central crop.  The same
two-colored strip argument applies to hexagonal cluster boundaries.
The collars absorb the square drawing's $O(\delta)$ displacement and
separate any shared interval endpoints before cropping.

A witnessing loop has spherical diameter bounded below by the
positive spherical distance between the two rims.  A sufficiently
close one-to-one matching therefore carries its witnessing intervals,
and their order and disjoint interiors, to matched loop circles.
It follows that $T_m(\,\cdot\,;a,b)$ is open and, when $4a<b$,
\begin{equation}\label{perc:even-closure}
 \overline{T_m(\,\cdot\,;a,b)}
 \subset T_m(\,\cdot\,;2a,b/2).
\end{equation}
For the closure assertion, match the witnessing loops to the finitely
many relevant limit loops and pass to subsequences of the interval
endpoints on their circles.  Cyclic order is preserved; the separated
rim visits and uniform continuity prevent a witnessing interval from
collapsing.  The enlarged rims make its limiting visits strict.

We also need local matching for triangular sites.  A monochromatic
circuit between a fixed window and a remote boundary confines every
loop meeting the window.  RSW makes the failure probability tend to
zero as the boundary recedes, uniformly in the mesh, so it suffices
first to work in a fixed disc.  In
\citet[Theorem~5]{CamiaNewman2006Full}, $K_\delta(\varepsilon)$ is a
tight bound on the steps needed to find every loop of diameter
$>\varepsilon$, and any fixed number of tracked exploration paths and
domain boundaries are jointly tight as ordered curves.  Each produced
loop is a bounded concatenation of subcurves of those tracked curves.
Compactness of subcurve endpoints gives tightness of the ordered loop
curves, while at most one loop per step gives tightness of their
number.  Diagonalizing over buffered diameter cutoffs therefore gives
subsequential tightness in one-to-one matching.

For any such matching limit, forgetting multiplicities and then roots
and orientations at positive diameter cutoffs gives the Hausdorff-set
limit of \citet[Theorem~1]{CamiaNewman2006Full}, which is the usual
$\CLE_6$ law \citep[Section~1.1]{Sheffield2009Exploration}.  There can be no
duplicate loop elements in the matching limit.  Every positive
$\CLE_6$ loop has an interior double visit with nonconstant
excursions, by Lemma~\ref{loop:continuation}.  Two copies would
therefore give at least three visits in total.  Matching convergence
would produce six disjoint radial medial sublegs around that point,
hence six alternating triangular arms.  The uniform six-arm bound of
\citet[Lemma~6.1]{CamiaNewman2006Full} excludes any such point as the
inner radius tends to zero, exactly as in
Lemma~\ref{loop:geometry}(ii).  Thus the local matching limit is
$\CLE_6$ with each loop counted once.  This argument uses the
ordered-curve metric for individual loops and does not infer
one-to-one matching from Hausdorff-set convergence alone.

Apply Portmanteau to the open event $T_m$ and its closure, using
\eqref{perc:even-lattice}--\eqref{perc:even-closure} and the common
$\CLE_6$ limit.  For $64r<R$ this gives the exact buffered comparison
\begin{equation}\label{perc:even-model-comparison}
 \limsup_{\delta\downarrow0}\pi^\square_{m,\delta}(r,R)
 \le
 \liminf_{\delta\downarrow0}\pi^\triangle_{m,\delta}(8r,R/8),
\end{equation}
and the same inequality with the lattices interchanged.  The
intermediate events for \eqref{perc:even-model-comparison} are
$T_m^\square(2r,R/2)$, its closure inside
$T_m^{\mathrm{CLE}}(4r,R/4)$, and the open event
$T_m^\triangle(4r,R/4)$.  Thus rim tangencies need not have
probability zero.

Let $g_m(t)$ be the triangular fixed-annulus limit for the alternating
word in the parallelogram annuli of \citet[Theorem~20]{Nolin2008NearCritical}.
That theorem and its following fixed-annulus statement, based on the
triangular-site arm calculation of \citet{SmirnovWerner2001Exponents},
give
$g_m(t)=t^{(m^2-1)/12+o(1)}$ as $t\downarrow0$.  Deterministic
inclusions between circular and parallelogram annuli change each rim
by a fixed factor.  Hence, for one fixed $K\ge64$ and each fixed
$0<t<1/K$,
\begin{equation}\label{perc:even-envelope}
 g_m(t/K)
 \le\liminf_{\delta\downarrow0}\pi^\square_{m,\delta}(t,1)
 \le\limsup_{\delta\downarrow0}\pi^\square_{m,\delta}(t,1)
 \le g_m(Kt).
\end{equation}
This is an envelope for fixed annuli, not an assertion that the
unbuffered square-bond annular probability has a limit.

\subsection{From fixed annuli to lattice exponents}

\begin{proof}[Completion of the proof of Corollary~\ref{intro:bond-exponents}]
At $q=1$, the arm separation, localization, and
quasi-multiplicativity results of
\citet[Propositions~6.2, 6.3, and 6.5]{DMT2021Fractal} apply to each
fixed square-bond color word.  Fixed-shape comparison transfers them
from their square annuli to the circular annuli used here.  In
particular, for $j=1$ or fixed even $j$ there are constants
$0<c_j\le C_j<\infty$ such that
\begin{equation}\label{perc:quasi-multiplicativity}
 c_j\pi_j(r,R)
 \le \pi_j(r,\rho)\pi_j(\rho,R)
 \le C_j\pi_j(r,R)
\end{equation}
at separated radii, where $\pi_j(r,R)$ denotes the corresponding
unit-mesh annular probability.  Equations
\eqref{perc:cle-one}, \eqref{perc:one-envelope}, and
\eqref{perc:even-envelope} say that, after taking the mesh limit
at each fixed ratio $M$, both envelopes have logarithmic cost
$\alpha_j\log M+o(\log M)$, with
$\alpha_1=5/48$ and $\alpha_j=(j^2-1)/12$ for even $j$.
Choose a large fixed $M$, then take all sufficiently large inner
radii so that their ratio-$M$ probabilities lie between the two
envelopes with an arbitrarily small logarithmic error.  Iterating
\eqref{perc:quasi-multiplicativity} over ratio-$M$ annuli adds at most
$\max\{|\log c_j|,|\log C_j|\}$ per annulus.  First let the number of
annuli tend to infinity, then divide by $\log M$ and let $M\to\infty$.
Monotonicity handles the last partial annulus,
and quasi-multiplicativity handles any fixed admissible inner radius.
A fixed finite inner connection compares the annular one-arm event
with connection from $0$.  This proves
\eqref{intro:bond-arm-exponents}.  The argument uses only fixed-ratio
limits followed by quasi-multiplicativity, not a uniform rate of
convergence in the annulus ratio.

\paragraph{Near-critical consequences.}
Let $L(p)$ be a finite-size crossing length with a fixed sufficiently
small RSW-admissible threshold, using the subcritical color on the
appropriate side of $p=1/2$.  Kesten's square-bond scaling relations
\citep{Kesten1987Scaling}, in the arm-event formulation discussed by
\citet[Sections~7--8.1]{Nolin2008NearCritical}, give
\begin{equation}\label{perc:scaling-relations}
\begin{gathered}
 |p-\tfrac12|L(p)^2\pi_4(r_2,L(p))\asymp1,\qquad
 \theta(p)\asymp\pi_1(L(p))\quad(p>\tfrac12),\\
 \chi(p)\asymp L(p)^2\pi_1(L(p))^2,\qquad
 \xi(p)\asymp L(p).
\end{gathered}
\end{equation}
Here $r_2$ is any fixed admissible inner radius for four arms; changing
that radius costs only a constant.  The second-moment length defined
with the Euclidean norm is comparable to the version with the
$\ell^\infty$ norm.  Substituting
$\pi_4(r_2,R)=R^{-5/4+o(1)}$ first gives
$L(p)=|p-\tfrac12|^{-4/3+o(1)}$.  Substituting
$\pi_1(R)=R^{-5/48+o(1)}$ in the remaining relations gives
$\beta=(4/3)(5/48)=5/36$ and
$\gamma=(4/3)(2-10/48)=43/18$, proving
\eqref{intro:bond-thermodynamic-exponents}.
\end{proof}
